% BHU PYQ -- Geography: Geography of India (2025-26 NEP)
\documentclass[11pt,a4paper]{article}
\usepackage[a4paper,top=2.5cm,bottom=2.5cm,left=2.8cm,right=2.8cm,headheight=14pt]{geometry}
\usepackage{amsmath,amssymb}
\usepackage{fontspec}
\setmainfont{Kohinoor Devanagari}
\usepackage{microtype,fancyhdr,enumitem,hyperref,lastpage,parskip}

\hypersetup{colorlinks=true,urlcolor=black,linkcolor=black,pdftitle={GGRMJ33: Geography of India -- BHU NEP 2025-26}}

\pagestyle{fancy}\fancyhf{}
\renewcommand{\headrulewidth}{0.4pt}\renewcommand{\footrulewidth}{0.4pt}
\fancyhead[L]{\small \textit{Banaras Hindu University}}
\fancyhead[R]{\small \textit{GGRMJ33: Geography of India (2025-26)}}
\fancyfoot[C]{\small Page \thepage\ of \pageref{LastPage}}

\newlist{parts}{enumerate}{2}
\setlist[parts,1]{label=(\alph*),leftmargin=*,itemsep=4pt,topsep=3pt}
\setlist[parts,2]{label=(\roman*),leftmargin=*,itemsep=2pt,topsep=2pt}

\newcommand{\pts}[1]{\hfill \textit{[#1]}}

\begin{document}

\begin{center}
  {\large\bfseries BANARAS HINDU UNIVERSITY}\\[2pt]
  {\normalsize B. A. / B. Sc. (Hons.) Semester III Examination, 2025-26 (NEP)}\\[6pt]
  \rule{0.8\linewidth}{0.5pt}\\[6pt]
  {\large\bfseries GEOGRAPHY}\\[4pt]
  {\large\bfseries Paper Code: GGRMJ33 : Geography of India}\\[4pt]
  {\normalsize Time: 2:30 Hours \hfill Full Marks: 70}\\[2pt]
  {\small REF NO. CECONF/N/2025-26/088}
\end{center}
\rule{\linewidth}{0.4pt}
\smallskip

\noindent \textit{Instructions: Answer five questions in all including Question No.\ 1 which is compulsory. Answer each part of Question No.\ 1 in approximately 200 words and remaining questions in approximately 1000 words. Illustrate your answers with suitable maps and diagrams. Write your Roll No.\ at the top immediately on receipt of this question paper.}

\smallskip
\noindent \textit{प्रश्न संख्या 1 (अनिवार्य) सहित कुल पाँच प्रश्नों के उत्तर दीजिए। प्रश्न 1 के प्रत्येक भाग का उत्तर लगभग 200 शब्दों में एवं शेष प्रश्नों का उत्तर लगभग 1000 शब्दों में दीजिए। अपने उत्तरों को उपयुक्त मानचित्रों तथा आरेखों द्वारा स्पष्ट कीजिए।}

\bigskip

\begin{enumerate}[label=\textbf{\arabic*.},leftmargin=*,itemsep=14pt]

  \item Write short notes on the following:\pts{4 $\times$ 3.5 = 14}
  
  निम्नलिखित पर संक्षिप्त टिप्पणियाँ लिखिए:
  \begin{parts}
    \item Geological structure of the Himalayan Mountains / हिमालय पर्वत की भूवैज्ञानिक संरचना
    \item Distribution of iron ore in India / भारत में लौह अयस्क का वितरण
    \item Major soil types of India / भारत की प्रमुख मृदा के प्रकार
    \item Population growth trends in India / भारत में जनसंख्या वृद्धि की प्रवृत्तियाँ
  \end{parts}

  \item Divide India into major physiographic divisions and describe the physical characteristics of the Northern Plains in detail.\pts{14}
  
  भारत को प्रमुख भौतिक/भू-आकृतिक प्रदेशों में विभाजित कीजिए तथा उत्तरी मैदान की भौतिक विशेषताओं का विस्तार से वर्णन कीजिए।

  \item Discuss the origin and mechanism of the Indian Monsoon. How does it affect Indian agriculture?\pts{14}
  
  भारतीय मानसून की उत्पत्ति एवं क्रियाविधि की विवेचना कीजिए। यह भारतीय कृषि को किस प्रकार प्रभावित करता है?

  \item Describe the spatial distribution and production of Wheat and Rice in India.\pts{14}
  
  भारत में गेहूं एवं चावल के स्थानिक वितरण तथा उत्पादन का वर्णन कीजिए।

  \item Describe the locational factors, production and distribution of Iron and Steel industry in India with special reference to Chota Nagpur region.\pts{14}
  
  छोटा नागपुर पठार क्षेत्र के विशेष संदर्भ में भारत में लौह एवं इस्पात उद्योग के अवस्थिति कारकों, उत्पादन एवं वितरण का वर्णन कीजिए।

  \item Divide India into major climate zones according to Koppen's scheme of classification.\pts{14}
  
  कोपेन के जलवायु वर्गीकरण की योजना के अनुसार भारत को प्रमुख जलवायु प्रदेशों में विभाजित कीजिए।

  \item Discuss the problems and prospects of industrial development in India.\pts{14}
  
  भारत में औद्योगिक विकास की समस्याओं एवं सम्भावनाओं की चर्चा कीजिए।

\end{enumerate}

\end{document}
